\section{Scope, boundaries and future interfaces}
\InputSectionBody{sections/11_discussion_conclusion}

\subsection{Exact boundary of the present article}

The reusable theorem in \cref{thm:one-well-composition} is an interface theorem: it composes exact sourced coefficients, an admissible observation/compression provider, actual-kernel local analysis, source-uniform weighted tail control and exact realization.  It does not assert that every finite local-map system has such a provider, and it does not identify the positive majorant with the actual complex kernel.

The moving-target conclusion in \cref{thm:compact-moving-target} is limited to one fixed reciprocal-semiprime family, fixed $t$, fixed admissible $\gamma$ and each fixed standardized compact window.  It does not give a growing-window or moderate-deviation theorem, does not exchange the family and target quantifiers, does not prove a joint value--cardinality local limit and does not imply a prescribed-family gap-free-floor theorem.

No mathematical appendix is needed for the proof.  In particular, this article does not import the full restricted-support or circuit branch, the maximal T1--T6 provider library, the thin-anchor exponent theory, the reciprocal-flow extension, or a second application.  It makes no coding, random-current, matroid-circuit, ordinary- or signed-flow, cycle-double-cover, AffineCDC, or prescribed-family consequence claim.  Those subjects require their own hypotheses and, where viable, their own mathematical home.

Within these boundaries the argument is complete: exact coefficient semantics lead to the reciprocal system; the one-anchor theorem supplies synchronization; rows and fibre compression supply the decoded tail; the actual kernel supplies the local Gaussian well; positivity, realization and no-wrap yield the direct target; prime dilution yields the characterization; and the same fixed family supplies compact local saturation and the stated quantitative consequences.
